\documentclass[11pt]{article}
\usepackage[T1]{fontenc}\usepackage{lmodern,amsmath,amssymb,booktabs}
\usepackage[margin=0.97in]{geometry}\usepackage[hidelinks]{hyperref}
\title{How much joint information certifies a Pareto cover after adaptive simulation?}
\author{Lluis Eriksson}\date{20 November 2026}
\begin{document}\maketitle
\begin{abstract}
A finite library of designs has known costs and expensive response vectors. Only a simulated design may serve as a witness. We ask how many evaluations guarantee that every design is covered at its cost and within specified additive tolerances. For a budget of \(k\) evaluations, joint support projections on at most \(k+1\) designs determine feasibility. Order \(k\) can fail even for binary scalar responses. Under an endpoint attainability condition the exact deterministic worst-case budget equals a static set-cover number. With dependence, two libraries agreeing on every \(t\)-design marginal distribution can require \(m\) and \(\min(m,t+2)\) evaluations respectively, while both fixed schedules require \(m\). The construction gives a limit on what low-order error summaries can certify.
\end{abstract}

\section{The simulation decision}
Braz\~ao and da Silva [1] compare neural and direct simulation for urban drainage. They evaluate proposed designs with the Storm Water Management Model, using peak flow, discharge volume and construction cost. Improvements in a front-distance metric for two storms motivate a separate question: how many evaluations guarantee coverage of an entire finite library if one physical design must serve across storms? We do not have the authors' design-level simulation data or a justified error envelope for that library.

Let \(I\) be the as-yet unevaluated designs, \(B\) an optional set of previously evaluated reference designs, \(c_i\) the known construction cost, and \(y_i\) a vector of losses to minimize. One query returns the entire vector at unit evaluation cost. For \(e\geq0\) coordinatewise, design \(r\) covers \(j\) if
\[
             c_r\leq c_j,\qquad y_{rh}\leq y_{jh}+e_h\quad\text{for every }h.
\]
For two storms stack their coordinates, with separate tolerances in physical units. Given known joint support \(\Omega\) of complete responses, a policy stops only if, in every scenario consistent with the observed values, every design is covered by some member of \(B\) or the queried set. Previously observed responses are fixed. Self-coverage is allowed. We study deterministic worst-case query count, with unrestricted processing of known support. Distinct minimizers may tie. A policy can always return all observed designs; discarding points exactly dominated by another observed design preserves coverage.

\newpage\section{A sharp projection order}
\textbf{Theorem 1.} For \(k\geq1\), two nonempty supports on the same labelled library, with the same reference values, costs and tolerances, admit the same valid policies of depth at most \(k\) if their projections on all subsets of at most \(k+1\) unevaluated designs agree. For every \(k\), projections on all \(k\)-subsets can agree while a budget \(k\) succeeds for only one support.

\emph{Proof.} Fix a policy and a transcript on queried labels \(Q\), \(|Q|\leq k\). Reachability is decided by the projection on \(Q\), since the policy uses only previous answers. Failure at a reachable leaf requires some unqueried label \(j\) with a value that violates coverage by every observed design. Whether that value and the transcript can coexist depends only on the projection on \(Q\cup\{j\}\). Conversely, a violating projected tuple extends to a feasible world producing the leaf. This proves the sufficient order and gives a leaf verification procedure. The statement also applies to infinite-valued supports when projections are actual attainable tuples.

For sharpness use \(k+1\) equal-cost designs, a reference of loss 2, and exact scalar coverage. Each response is 0 or 2. Compare the full binary cube with its subset in which the number of zero losses is even. The uniform laws have the same marginals on \(k\) labels: every assignment leaves one bit uniquely determined by parity in the even model. Query any \(k\) designs. An observed zero covers everyone. If every answer is 2, the last loss must also be 2 under even parity. In the full cube the same all-high transcript can hide a single zero at any unqueried design. Under even parity fewer than \(k\) all-high answers can hide two zeros. The respective optimal budgets are exactly \(k+1\) and \(k\). \(\square\)

\paragraph{When adaptive choice has zero worst-case gain.}
Suppose every unknown \(y_i\) lies between finite coordinatewise bounds \(L_i,U_i\). Assume the all-upper response is feasible and, for every \(j\), the response that replaces only \(U_j\) by \(L_j\) is feasible. Form a directed cover graph: an observed reference covers \(j\) if its value covers \(L_j\); an unknown \(r\ne j\) has an edge to \(j\) if \(c_r\leq c_j\) and \(U_r\leq L_j+e\). Give each queried design its own label as well.

\textbf{Theorem 2.} The minimum number of selected vertices covering all labels, including the free reference coverage, equals both the optimal fixed budget and the optimal adaptive worst-case budget.

\emph{Proof.} A graph cover works for all responses between the bounds. Follow any correct adaptive policy on the all-upper world until it stops at \(Q\). If \(j\notin Q\) has no covering edge from the reference or \(Q\), the single replacement at \(j\) preserves the transcript and defeats every observed witness. Thus \(Q\) is a graph cover. The same fixed \(Q\) suffices. No probabilistic independence or full box of attainable values was used. \(\square\)

\newpage\section{Identical low-order laws, different optimal budgets}
\textbf{Theorem 3.} For integers \(m,t\geq1\), there are two finite laws on the same \(m+t+1\) labels with identical marginals on every set of at most \(t\) labels. The optimal fixed budget is \(m\) for both. Their adaptive worst-case budgets are \(m\) and \(\min(m,t+2)\).

\emph{Construction.} All construction costs are equal. A reference has loss 2. The \(m\) candidates independently have loss 0 or 2 with equal probability. Define \(s\) as the first zero's index in \(1,\ldots,m\), or 0 if no zero exists. The other \(t+1\) designs have losses \(3+d_i\), with digits in the additive group modulo \(m+1\). They are already covered by the reference but their evaluations may carry information. In the product law all digits are independent uniform and independent of the candidates. In the coupled law draw \(t\) digits independently and uniformly, and set the last so that \(\sum_i d_i=s\) modulo \(m+1\). This is an elementary use of additive sharing [6].

\emph{Marginal equality.} Conditional on the entire candidate vector, prescribing \(q\leq t\) probe digits leaves \((m+1)^{t-q}\) tuples with the required sum. Hence those digits have the uniform product law independent of every candidate response. This proves the stated equality of distributions, not just ranges.

\emph{Upper and lower bounds.} Any fixed schedule omitting candidate \(j\) fails when \(j\) is the unique zero. Reading all \(m\) candidates suffices. In the product model a transcript with all candidates observed so far high can always hide a zero at an unread label, for any digits. Its adaptive worst-case cost is therefore \(m\). In the coupled model one can query all digits, and, when their sum is nonzero, query the indicated candidate to obtain the required witness. This gives the upper bound.

For the lower bound consider any proposed budget \(d<\min(m,t+2)\). Answer a candidate query with 2 and each of the first \(t\) distinct probe queries with digit 0. While one probe remains unread and fewer than \(m\) candidates have been queried, the transcript is compatible with a unique zero at an unread candidate: use a missing digit to complete the sum. The reference fails to cover that zero. If all \(t+1\) probes are queried within the proposed budget, there was no room for any candidate query. On the final probe choose the digit giving sum 1. The policy has identified candidate 1 but has never evaluated it, so cannot supply the witness. These cases exclude every smaller budget. Repeated queries add no information. \(\square\)

\paragraph{Two storms.} Each scalar value may be used in all four response coordinates, preserving the comparison and making responses nondecreasing with severity. This embedding has redundant objectives. It demonstrates a limitation of marginal validation even in a simple subclass. The informative designs are synthetic. There is no claim that hydraulic simulators encode a modular signal.

\newpage\section{Interpretation, checks and scope}
Under Theorem 2 one can quote the graph optimum as a worst-case cost in simulation units. Theorem 1 gives a different route if joint support information is available. The smaller projections can falsify unsafe stopping rules without specifying a full probabilistic model, provided the necessary tuple information is available. Enumerating or estimating such projections may itself be costly. In particular, estimates of pairwise correlation are weaker than exact pairwise distributions. Theorem 3 shows that even exact distributions of fixed finite order can give arbitrarily different answers about the benefit of adaptivity as the library grows.

For example, for five candidates and three probes, the two laws agree on all pairs. Both fixed schedules need five evaluations. Four evaluations suffice and are necessary in the coupled model. The supplement computes this from all \(2^5 6^3\) product worlds and the \(2^5 6^2\) coupled worlds. It minimizes over the next query at each possible transcript, checks the projected probabilities with integer arithmetic, and compares the graph test against complete scenarios on a two-objective, unequal-cost instance. These finite checks supplement the proofs.

\paragraph{Relation to earlier work.} One-exact approximate Pareto sets [2] retain one exact objective. Pareto active learning [3] studies statistical identification with confidence bounds. Reconstruction of imprecise Pareto fronts [4] studies retrievals and instance-sensitive geometric algorithms. Interactive submodular cover [5] already combines learning and covering. We claim neither the general idea of informative probes nor a new sharing scheme. Our conclusions concern the projection order for a budget with observed witnesses, an endpoint condition for exact fixed scheduling, and matching costs for two low-order-indistinguishable libraries. We do not assert priority across all decision-tree formulations. The results do not estimate the cost saving for a particular catchment, expected performance, or a hydraulic uncertainty distribution.

\paragraph{Model assessments and review context.} OpenAI Codex participated in the derivation, tests and writing. The author administers AIRR. The version-specific assessment is attached separately. The simulations in the motivating preprint have not been independently reproduced here.

\begin{thebibliography}{6}\small
\bibitem{water} A. J. da Costa Braz\~ao and G. B. Lima da Silva. Assessing the reliability of ANN-assisted multiobjective optimization for urban low impact development design. Research Square, v1 (2026). \href{https://doi.org/10.21203/rs.3.rs-11446856/v1}{Preprint and data statement}.
\bibitem{exact} A. Herzel et al. One-exact approximate Pareto sets. J. Global Optimization 80 (2021), 87--115. \href{https://doi.org/10.1007/s10898-020-00951-7}{Article}.
\bibitem{pal} M. Zuluaga et al. Active learning for multi-objective optimization. ICML (2013). \href{https://proceedings.mlr.press/v28/zuluaga13.html}{Proceedings}.
\bibitem{front} S. de Berg et al. Instance optimal and universally optimal bounds for imprecise Pareto fronts. ESA (2026). \href{https://doi.org/10.4230/LIPIcs.ESA.2026.106}{Proceedings}.
\bibitem{cover} A. Guillory and J. Bilmes. Interactive submodular set cover (2010). \href{https://arxiv.org/abs/1002.3345}{Authors' preprint}.
\bibitem{share} A. R. Ghavamipour et al. Privacy-preserving logistic regression with secret sharing (2022). \href{https://doi.org/10.1186/s12911-022-01811-y}{Additive sharing construction}.
\end{thebibliography}
\end{document}

